\ParallelTitle{Keyword Quick Reference}{关键词速查表}
\ParallelText{quick-intro}
 {Find a remembered word or task in the A--Z lists, then follow the canonical entry and its Notes link. Every listed English headword has its own row. Chinese terms mirror that order. Page numbers are generated automatically.}
 {在 A--Z 列表中查找记得的词语或任务，再打开概念条目及其“笔记”链接。每个英文检索词都有独立的一行，中文与之对应。页码自动生成。}
\ParallelSection{aliases}{Aliases and acronyms}{别称与缩写}
\ParallelText{alias-last-in-first-out}
 {\textbf{Last in, first out.} See \ParallelReference{concept-stack}{Stack behaviour}.}
 {\textbf{后进先出。}见\hyperref[concept-stack]{栈的行为（第\pageref*{concept-stack}页）}。}
\ParallelText{alias-lifo}
 {\textbf{LIFO.} See \ParallelReference{concept-stack}{Stack behaviour}.}
 {\textbf{LIFO。}见\hyperref[concept-stack]{栈的行为（第\pageref*{concept-stack}页）}。}
\ParallelText{alias-mean}
 {\textbf{Mean speed.} See \ParallelReference{concept-speed}{Average speed}.}
 {\textbf{Mean speed（平均速率的英文别称）。}见\hyperref[concept-speed]{平均速率（第\pageref*{concept-speed}页）}。}
\ParallelSection{keywords}{Keywords and tasks}{关键词与任务}
\ParallelText{key-area}
 {\textbf{Area.} See \ParallelReference{concept-area}{Rectangle area}.}
 {\textbf{面积。}见\hyperref[concept-area]{矩形面积（第\pageref*{concept-area}页）}。}
\ParallelText{key-capacity}
 {\textbf{Capacity.} See \ParallelReference{concept-stack}{Stack behaviour}.}
 {\textbf{容量。}见\hyperref[concept-stack]{栈的行为（第\pageref*{concept-stack}页）}。}
\ParallelText{key-cover}
 {\textbf{Cover a surface.} See \ParallelReference{concept-area}{Rectangle area}.}
 {\textbf{铺满表面。}见\hyperref[concept-area]{矩形面积（第\pageref*{concept-area}页）}。}
\ParallelText{key-dimensions}
 {\textbf{Dimensions.} See \ParallelReference{concept-area}{Rectangle area}.}
 {\textbf{尺寸。}见\hyperref[concept-area]{矩形面积（第\pageref*{concept-area}页）}。}
\ParallelText{key-displacement}
 {\textbf{Displacement.} See \ParallelReference{concept-speed}{Average speed}.}
 {\textbf{位移。}见\hyperref[concept-speed]{平均速率（第\pageref*{concept-speed}页）}。}
\ParallelText{key-distance}
 {\textbf{Distance.} See \ParallelReference{concept-speed}{Average speed}.}
 {\textbf{路程。}见\hyperref[concept-speed]{平均速率（第\pageref*{concept-speed}页）}。}
\ParallelText{key-doubling}
 {\textbf{Doubling.} See \ParallelReference{concept-area}{Rectangle area}.}
 {\textbf{加倍。}见\hyperref[concept-area]{矩形面积（第\pageref*{concept-area}页）}。}
\ParallelText{key-edging}
 {\textbf{Edging.} See \ParallelReference{concept-area}{Rectangle area}.}
 {\textbf{沿边铺设。}见\hyperref[concept-area]{矩形面积（第\pageref*{concept-area}页）}。}
\ParallelText{key-elapsed}
 {\textbf{Elapsed time.} See \ParallelReference{concept-speed}{Average speed}.}
 {\textbf{总时间。}见\hyperref[concept-speed]{平均速率（第\pageref*{concept-speed}页）}。}
\ParallelText{key-empty-collection}
 {\textbf{Empty collection.} See \ParallelReference{concept-stack}{Stack behaviour}.}
 {\textbf{空集合。}见\hyperref[concept-stack]{栈的行为（第\pageref*{concept-stack}页）}。}
\ParallelText{key-empty-stack}
 {\textbf{Empty stack.} See \ParallelReference{concept-stack}{Stack behaviour}.}
 {\textbf{空栈。}见\hyperref[concept-stack]{栈的行为（第\pageref*{concept-stack}页）}。}
\ParallelText{key-fifo}
 {\textbf{FIFO.} Compare with \ParallelReference{concept-stack}{Stack behaviour}.}
 {\textbf{FIFO。}与\hyperref[concept-stack]{栈的行为（第\pageref*{concept-stack}页）}比较。}
\ParallelText{key-overflow}
 {\textbf{Overflow.} See \ParallelReference{concept-stack}{Stack behaviour}.}
 {\textbf{溢出。}见\hyperref[concept-stack]{栈的行为（第\pageref*{concept-stack}页）}。}
\ParallelText{key-pauses}
 {\textbf{Pauses.} See \ParallelReference{concept-speed}{Average speed}.}
 {\textbf{停顿。}见\hyperref[concept-speed]{平均速率（第\pageref*{concept-speed}页）}。}
\ParallelText{key-perimeter}
 {\textbf{Perimeter.} See \ParallelReference{concept-area}{Rectangle area}.}
 {\textbf{周长。}见\hyperref[concept-area]{矩形面积（第\pageref*{concept-area}页）}。}
\ParallelText{key-pop}
 {\textbf{Pop.} See \ParallelReference{concept-stack}{Stack behaviour}.}
 {\textbf{出栈。}见\hyperref[concept-stack]{栈的行为（第\pageref*{concept-stack}页）}。}
\ParallelText{key-push}
 {\textbf{Push.} See \ParallelReference{concept-stack}{Stack behaviour}.}
 {\textbf{入栈。}见\hyperref[concept-stack]{栈的行为（第\pageref*{concept-stack}页）}。}
\ParallelText{key-queue}
 {\textbf{Queue.} Compare with \ParallelReference{concept-stack}{Stack behaviour}.}
 {\textbf{队列。}与\hyperref[concept-stack]{栈的行为（第\pageref*{concept-stack}页）}比较。}
\ParallelText{key-removal-order}
 {\textbf{Removal order.} See \ParallelReference{concept-stack}{Stack behaviour}.}
 {\textbf{移除顺序。}见\hyperref[concept-stack]{栈的行为（第\pageref*{concept-stack}页）}。}
\ParallelText{key-scaling}
 {\textbf{Scaling.} See \ParallelReference{concept-area}{Rectangle area}.}
 {\textbf{缩放。}见\hyperref[concept-area]{矩形面积（第\pageref*{concept-area}页）}。}
\ParallelText{key-uncertainty}
 {\textbf{Uncertainty.} See the gap under \ParallelReference{concept-speed}{Average speed}.}
 {\textbf{不确定度。}见\hyperref[concept-speed]{平均速率（第\pageref*{concept-speed}页）}中的缺口。}
\ParallelText{key-unit-square}
 {\textbf{Unit square.} See \ParallelReference{concept-area}{Rectangle area}.}
 {\textbf{单位正方形。}见\hyperref[concept-area]{矩形面积（第\pageref*{concept-area}页）}。}
\ParallelText{key-units}
 {\textbf{Units.} Rates: \ParallelReference{concept-speed}{Average speed}. Squared units: \ParallelReference{concept-area}{Rectangle area}.}
 {\textbf{单位。}速率单位见\hyperref[concept-speed]{平均速率（第\pageref*{concept-speed}页）}；平方单位见\hyperref[concept-area]{矩形面积（第\pageref*{concept-area}页）}。}
\clearpage
\ParallelSection{concept-speed}{Average speed}{平均速率}
\ParallelText{recall-speed}
 {\textbf{Rule:} $v_{\mathrm{avg}}=d/\Delta t$, with total path length $d$ and elapsed time $\Delta t>0$, including pauses. Metres divided by seconds gives m/s. Example: $(60+120)/12=15$ m/s.}
 {\textbf{规则：}$v_{\mathrm{avg}}=d/\Delta t$，其中 $d$ 为总路程，$\Delta t>0$ 为总时间，包含停顿。米除以秒，得到 m/s。例如：$(60+120)/12=15$ m/s。}
\ParallelText{recall-speed-limit}
 {\textbf{Remember:} distance is not displacement. Two unequal speeds can be averaged arithmetically only for equal durations. This average specifies neither instantaneous nor maximum speed.}
 {\textbf{注意：}路程不是位移。两个不同速率只有在用时相等时，才能用算术平均值得到全程平均速率。平均速率不能确定瞬时速率或最高速率。}
\ParallelText{recall-speed-check}
 {\textbf{Check:} multiply the answer by elapsed time to recover distance. This is an arithmetic check. \textbf{Gap:} the source gives no measurement-uncertainty method.}
 {\textbf{核验：}将结果乘以总时间，反算路程。这是计算核验。\textbf{缺口：}原材料没有给出测量不确定度的估计方法。}
\ParallelText{qr-notes-link-speed}
 {\hyperref[notes-speed]{Notes, section~\ref*{notes-speed}, p.~\pageref*{notes-speed}}}
 {\hyperref[notes-speed]{笔记第\ref*{notes-speed}节，第\pageref*{notes-speed}页}}
\ParallelSection{concept-area}{Rectangle area}{矩形面积}
\ParallelText{recall-area}
 {\textbf{Rule:} $A=wh$ for positive rectangle side lengths in the same length unit; the area unit is squared. Convert first: 200 cm = 2 m, so $3\times2=6\ \mathrm{m^2}$. Non-integer side lengths are allowed.}
 {\textbf{规则：}矩形的 $A=wh$，两个边长均为正数，且使用同一长度单位；面积单位是该单位的平方。先换算：200 cm = 2 m，所以 $3\times2=6\ \mathrm{m^2}$。边长可以不是整数。}
\ParallelText{recall-area-scale}
 {\textbf{Scaling:} multiply both sides by $k>0$ and area becomes $k^2A$. Double one side: $2A$; double both: $4A$.}
 {\textbf{缩放：}两个边长都乘以 $k>0$，面积变为 $k^2A$。只把一边加倍：$2A$；两边都加倍：$4A$。}
\ParallelText{recall-area-choice}
 {\textbf{Choose:} covering uses area; edging uses perimeter $P=2(w+h)$. For 3 m by 2 m, $A=6\ \mathrm{m^2}$ but $P=10$ m. Two adjacent lengths alone do not determine an arbitrary quadrilateral's area.}
 {\textbf{选择：}铺满表面用面积；沿边铺设用周长 $P=2(w+h)$。3 m 乘 2 m 的矩形，$A=6\ \mathrm{m^2}$，而 $P=10$ m。仅凭两条相邻边长，不能确定任意四边形的面积。}
\ParallelText{qr-notes-link-area}
 {\hyperref[notes-area]{Notes, section~\ref*{notes-area}, p.~\pageref*{notes-area}}}
 {\hyperref[notes-area]{笔记第\ref*{notes-area}节，第\pageref*{notes-area}页}}
\ParallelSection{concept-stack}{Stack behaviour}{栈的行为}
\ParallelText{recall-stack}
 {\textbf{Rule:} last in, first out (LIFO). Push adds at the top; pop removes and returns the top item. Write bottom to top, top at right. Behaviour does not dictate array or linked storage.}
 {\textbf{规则：}后进先出（LIFO）。入栈在栈顶添加元素；出栈移除并返回栈顶元素。这里从左到右按栈底到栈顶书写。行为规则不限定数组或链式存储。}
\ParallelText{recall-stack-example}
 {\textbf{Trace:} $[\ ]\xrightarrow{\text{push }A}[A]\xrightarrow{\text{push }B}[A,B]$. Pop returns $B$ and leaves $[A]$; the next returns $A$ and leaves $[\ ]$. FIFO would return $A$ first.}
 {\textbf{过程：}$[\ ]\xrightarrow{\text{入栈 }A}[A]\xrightarrow{\text{入栈 }B}[A,B]$。出栈返回 $B$，留下 $[A]$；再出栈返回 $A$，留下 $[\ ]$。FIFO 则先返回 $A$。}
\ParallelText{recall-stack-check}
 {\textbf{Test:} returned values and remaining state, empty-pop handling, and overflow if capacity is fixed. \textbf{Unspecified:} empty/overflow responses and implementation. Follow the actual interface contract.}
 {\textbf{测试：}返回值、剩余状态、空栈出栈的处理，以及容量固定时的溢出处理。\textbf{未规定：}空栈与溢出的响应，以及实现方式。应遵循实际接口约定。}
\ParallelText{qr-notes-link-stack}
 {\hyperref[notes-stack]{Notes, section~\ref*{notes-stack}, p.~\pageref*{notes-stack}}}
 {\hyperref[notes-stack]{笔记第\ref*{notes-stack}节，第\pageref*{notes-stack}页}}
